% !TeX program = lualatex
% Compile twice: lualatex 172.tex
% All references and prose are embedded; no BibTeX database or external inputs.
\documentclass[11pt,a4paper]{article}
\usepackage[margin=24mm,headheight=14pt]{geometry}
\usepackage{fontspec}
\setmainfont{Latin Modern Roman}
\setsansfont{Latin Modern Sans}
\setmonofont{Latin Modern Mono}
\usepackage{amsmath,amssymb,mathtools}
\usepackage{unicode-math}
\setmathfont{Latin Modern Math}
\usepackage{microtype}
\usepackage{enumitem,xurl,needspace}
\usepackage[dvipsnames]{xcolor}
\usepackage{hyperref}
\hypersetup{unicode=true,colorlinks=true,linkcolor=MidnightBlue,urlcolor=MidnightBlue,
 pdftitle={500 Mathematical Problem Statements},pdfauthor={Source-annotated compilation},bookmarksnumbered=true}
\usepackage{bookmark}
\usepackage{tocloft}
\setlength{\cftsecnumwidth}{3em}
\cftsetpnumwidth{2.5em}
\cftsetrmarg{4em}
\usepackage{titlesec}
\titleformat{\section}{\Large\bfseries\raggedright}{\thesection}{0.7em}{}
\usepackage{fancyhdr}
\pagestyle{fancy}\fancyhf{}
\fancyhead[L]{\small\sffamily Mathematical problem statements}
\fancyhead[R]{\small Source edition 22}
\fancyfoot[C]{\thepage}
\setlength{\parindent}{0pt}
\setlength{\parskip}{5pt plus 1pt}
\setlength{\emergencystretch}{3em}
\setcounter{tocdepth}{1}
\setcounter{secnumdepth}{1}
\setlist[itemize]{leftmargin=3em,itemsep=5pt,topsep=4pt,parsep=0pt}
\allowdisplaybreaks
\newcommand{\N}{\mathbb{N}}
\newcommand{\Z}{\mathbb{Z}}
\newcommand{\Q}{\mathbb{Q}}
\newcommand{\R}{\mathbb{R}}
\newcommand{\C}{\mathbb{C}}
\newcommand{\F}{\mathbb{F}}
\newcommand{\A}{\mathbb{A}}
\newcommand{\PP}{\mathbb P}
\newcommand{\E}{\mathbb{E}}
\newcommand{\eps}{\varepsilon}
\newcommand{\dd}{\,\mathrm d}
\newcommand{\sref}[2]{\hyperlink{source:#1:#2}{[S#2]}}
\newcommand{\eref}[2]{\hyperlink{extra:#1:#2}{[E#2]}}
% Turn the next switch off for a shorter reading copy. The quotations preserve
% every exactTarget field, including qualifications omitted from a short summary.
\newif\ifshowcatalogue
\showcataloguetrue
\newcommand{\cataloguescope}[1]{\ifshowcatalogue
 \paragraph{Exact scope in the supplied catalogue.}
 \begin{quote}\small #1\end{quote}\fi}

% Standalone notebook wrapper; original problem section retained verbatim.
\numberwithin{equation}{section}
\hypersetup{pdftitle={Research attempt -- problem 172},
 pdfauthor={Research notebook; source compilation retained}}
\fancyhead[L]{\small\sffamily Research notebook}
\fancyhead[R]{\small\sffamily Problem 172 | 27 September 2026}
\begin{document}
\setcounter{section}{171}
% ================================================================
% problemId: problem.separation-of-p-r-from-np-r-in-the-blum-shub-smale-model
% source releaseRank: 172; source releaseStatus: open
% exactTarget SHA-256: 5bbb68ca760a62d98876f6c0f89bb49e0531b6f12557aba1538ff786571b4085
\clearpage
\section{\texorpdfstring{Separation of P\_R from NP\_R in the Blum–Shub–Smale Model}{Separation of P\_R from NP\_R in the Blum–Shub–Smale Model}}\label{172}
\subsection{Definitions and mathematical statement}
An ordered-real Blum--Shub--Smale machine stores exact real numbers and performs field operations and equality/order tests at unit cost, using finitely many built-in real constants. Input length is the number of real coordinates, not their binary precision. Let $\mathrm P_{\R}$ consist of sets decided in polynomially many deterministic steps. A set belongs to $\mathrm{NP}_{\R}$ if membership is equivalent to existence of a polynomial-length real witness accepted by such a verifier. The assertion is
\[
 \mathrm P_{\R}\ne\mathrm{NP}_{\R}.
\]
The real constants belong to the fixed machine; it does not receive arbitrary input-dependent advice. This is a real-computation separation, not merely the classical binary $\mathrm P$ versus $\mathrm{NP}$ question.
\par\smallskip\textit{Formulation and concept references:} \sref{172}{1}.
\cataloguescope{In the standard ordered-real Blum–Shub–Smale model, machines use finitely many built-in real constants, real-valued registers, exact field operations (+,-,multiplication,division where defined) and order/equality tests at unit cost. Inputs are finite real vectors and size is vector length. Let P\_R be decision problems solved by deterministic machines in polynomially many steps, and NP\_R those with polynomial-length real witnesses verifiable in P\_R. The assertion is P\_R != NP\_R.}
\subsection{Short English statement}
In exact real-number computation, are some problems with short real certificates impossible to solve in polynomially many deterministic operations?
% BEGIN RESEARCH ADDITION ATTEMPT-172
\subsection{Research attempt}

\paragraph{Current frontier.}
The ordered-real model and its polynomial feasibility problems are described
in \sref{172}{1}; its comparisons with binary computation do not identify
the two models. In particular, a large binary encoding of a register is
not a lower bound on the number of BSS steps. The source explains the
complete-problem and geometric viewpoints, but does not establish the
separation requested here. No theorem verified in this research establishes
that separation with the finite arbitrary-real-constant convention intact.

\paragraph{Attack route.}
One route seeks lower bounds from the topology of semialgebraic decision
sets. Another tries to transfer binary hardness, but must account for exact
arithmetic and built-in real constants. A third would diagonalize against
real machines; a countable enumeration of rational-constant programs does
not cover the present machines. We pursue a concrete topological counting
argument and challenge it with a family whose topology is exponentially
complicated but whose computation is short. This tests the quantitative
strength of the chosen invariant before using it on a complete problem.

\paragraph{Attempt: degree growth along a computation.}
Fix a machine, an input dimension, and a line in the input space, with its
coordinates affine functions of a real parameter $x$. Its built-in real
constants are treated as coefficients, not as advice chosen separately for
each input. On any fixed computation path of at most $T$ operations, every
register is a rational function of $x$ whose numerator and denominator
have degree at most $2^T$.

To prove the bound, start with affine polynomials and constant denominators.
If the degrees at one stage are at most $D$, addition or subtraction of two
fractions gives numerator $ad\pm bc$ and denominator $bd$, of degrees at
most $2D$. Multiplication and division have the same bound, when the latter
is defined. Cancellation can only lower the degrees. Induction proves the
claim. No bound on coefficient sizes, and no algebraicity of the real
constants, was assumed.

An order or equality test of a rational expression can change its outcome
only at a zero of its numerator or denominator. An identically zero
polynomial gives a fixed test outcome and contributes no roots. Along one
path there are at most $T$ tests and division-domain conditions, so their
exceptional points number at most $C T2^T$. There are at most $3^T$ test
paths, since a sign test has at most three outcomes. After collecting all
such roots, every remaining open interval follows a fixed path and has a
fixed acceptance outcome. Including the isolated exceptional points gives
\begin{equation}\label{eq172-components}
 \#\{\text{connected components on the input line}\}
                  \le C'(T+1)6^T.
\end{equation}
The constants absorb harmless conventions about arithmetic nodes and halt
states. At exceptional points the machine is still required to decide
correctly, but they can add only their number of singleton components.

Thus a line section with $K$ connected components yields only
$T\ge c\log K-O(\log\log K)$ by this argument. Exponentially many
components in a dimension parameter supply a linear or polynomial bound,
not automatically a superpolynomial one. This does not rule out every
possible topological lower-bound method; it identifies the limitation of
this particular one-dimensional counting statistic.

\paragraph{Attempt: an exact fast family with exponential topology.}
Define
\[
 P_0(x)=x,\qquad P_{j+1}(x)=2P_j(x)^2-1.
\]
The cosine identity $2\cos^2\theta-1=\cos(2\theta)$ proves by induction
\begin{equation}\label{eq172-chebyshev}
              P_j(\cos\theta)=\cos(2^j\theta).
\end{equation}
For $j\ge1$, the set
$\{x\in[-1,1]:P_j(x)>0\}$ has exactly $2^{j-1}+1$ connected components.
Indeed $\theta\mapsto\cos\theta$ is a bijection from $[0,\pi]$ to
$[-1,1]$, and the $2^j$ simple zeros of $\cos(2^j\theta)$ alternate its
sign, with positive intervals at both endpoints. Nevertheless evaluation
uses only $j$ squarings, $j$ multiplications by $2$, $j$ subtractions, and
one sign test.

To make this a legitimate variable-size BSS decision family, let an input
have $j+1$ real coordinates, use the first as $x$, and use the remaining
coordinates only to specify length. A fixed uniform machine loops $j$
times and decides membership in polynomial, in fact linear, time.
Padding does not give input-dependent constants. Hence exponentially many
components and degree $2^j$ can coexist with membership in $\mathrm P_{\R}$.
This exact example falsifies the proposed inference from exponential
algebraic complexity to superpolynomial running time.

It also has a short real-witness presentation: introduce $y_0,\ldots,y_j$
and impose
\[
 y_0=x,\qquad y_{i+1}=2y_i^2-1,\qquad y_j>0.
\]
These are quadratic equations and one inequality. They illustrate why
introducing witnesses for arithmetic gates does not in itself create a
hard instance: here the witnesses are uniquely determined and quickly
computed. Real feasibility as a complete problem, discussed in
\sref{172}{1}, requires hardness for arbitrary systems, not every special
system produced by a short straight-line recurrence.

\paragraph{Attempt: where a stronger strategy would have to change.}
One might try to prove that every accepting computation for a carefully
chosen feasibility instance requires high-degree tests. The degree bound
above would then turn a degree lower bound $D$ into only $T\ge\log_2D$.
A superpolynomial time conclusion from that route would require a
correspondingly enormous degree lower bound, or a complexity measure which
cannot be cheaply generated by repeated squaring. The Chebyshev family
shows why merely counting real roots of a high-degree polynomial is not
that measure.

The binary-hardness route has a different obstruction. On binary inputs,
a real machine can generate integers with exponentially many binary digits
in polynomially many squarings. A faithful Turing simulation must pay for
such arithmetic; assuming its cost is polynomial would silently change
models. Conversely, a single fixed real constant cannot simply be used as
a unit-cost table lookup of arbitrary bits: a valid BSS algorithm still
must extract any purported information with its allowed field operations
and tests. Neither side of that informal argument yields a separation.

Finally, enumerating only machines whose constants are rational or algebraic
would prove, at most, a different statement. The path-degree calculation
is robust under arbitrary real coefficients, but a diagonal construction
would need a similarly robust treatment of all choices of finitely many
real constants. No such diagonal language with the required
$\mathrm{NP}_{\R}$ verification has been constructed here.

\paragraph{Stress test.}
The component bound counts the whole decision tree, not just the path of
one sample input. Denominator roots are retained, so division has not been
treated as a globally defined polynomial operation. The bound is on a line
section, not on arbitrary high-dimensional Betti numbers. A negative result
for restricted algebraic computation trees would also require justification
before being transferred to the full machine model.

The polynomial family above uses only the constants $1$ and $2$, and its
component count follows from an exact trigonometric identity. Numerical
plots or counting approximate zeros are unnecessary. It is not a
counterexample to $\mathrm P_{\R}\ne\mathrm{NP}_{\R}$: it is an easy
language deliberately chosen to refute a candidate lower-bound lemma.
No implication to binary $\mathrm P\ne\mathrm{NP}$ or its converse has
been presumed.

\paragraph{Outcome.}
\textbf{Outcome: Exploratory approach with unresolved gap.}

The attempt proves a path-degree bound and a corresponding component-count
bound valid with arbitrary fixed real constants. It then gives an explicit
family showing that exponential degree and topology are insufficient for
the proposed separation argument. These elementary deductions identify a
failed route, not a new lower bound for a complete real-feasibility problem.
A successful continuation needs a stronger invariant or a different
uniform hardness argument that survives exact arithmetic and the full
constant convention.
% END RESEARCH ADDITION ATTEMPT-172
\subsection{Sources}
\begin{itemize}\raggedright
\item[\textnormal{[S1]}]\hypertarget{source:172:1}{} Computing over the Reals: Where Turing Meets Newton, author exposition associated with2004 publication.
\url{https://www.cs.cmu.edu/~lblum/PAPERS/TuringMeetsNewton.pdf}.
{\small\textit{Catalogue formulation source.}}
\end{itemize}

\end{document}
